% CN 10.12 Exercises (p.385--388)

% ======================================================================
\begin{exercise}{1}
    Two independent fair dice are thrown, represented by random variables $X$
    and $Y$. Calculate the probability distribution of $|X-Y|$.
\end{exercise}

\begin{exercise}{2}
    Suppose $X$ and $Y$ are independent $\{0,1\}$-valued random variables where
    $X$ is unbiased, $\Pr(X=0)=\Pr(X=1)=\frac12$, but the distribution of $Y$ is
    unknown. What can you say about the probability distribution of
    $X\oplus Y$ (i.e.\ $X+Y\bmod2$)?
\end{exercise}

\begin{exercise}{3}
    One of the bitstrings $000,011,101,110$ is chosen uniformly at random. Let
    $X,Y,Z$ be its first, second and third bits.
    \begin{enumerate}[label=(\alph*)]
        \item Prove that (i) $X$ and $Y$ are independent; (ii) $X$ and $Z$ are
              independent; (iii) $Y$ and $Z$ are independent.
        \item Determine whether $Z$ is independent of the pair $(X,Y)$.
        \item Similarly, determine whether $Y$ is independent of the other two,
              and whether $X$ is independent of the other two.
    \end{enumerate}
\end{exercise}

\begin{exercise}{4}
    You need a die but don't have one. Your friend suggests tossing five fair
    coins, counting the Heads (in $\{0,\dots,5\}$) and adding one to get a number
    in $\{1,\dots,6\}$. Discuss this suggestion. Can you think of any other way
    in which a set of fair coins (possibly more than six) can be used to simulate
    the throw of a fair die?
\end{exercise}

\begin{exercise}{5}
    Define the random variable $W$ by
    \begin{center}
    \begin{tabular}{c|c}
        $k$ & $\Pr(W=k)$\\\hline
        0 & $\frac12$\\
        1 & $\frac13$\\
        2 & $\frac16$
    \end{tabular}
    \end{center}
    Determine the expectation, median, mode, variance and standard deviation of
    $W$.
\end{exercise}

\begin{exercise}{6}
    Find the variance of the sum of the numbers shown on two fair dice. What is
    the probability that this sum is at least two standard deviations away from
    its mean? Compare with what Chebyshev's Inequality says.
\end{exercise}

\begin{exercise}{7}
    Suppose $X\sim\mathrm{Bin}(16,0.5)$.
    \begin{enumerate}[label=(\alph*)]
        \item Your friend claims that, because $p=\frac12$, half the trials must
              be successes and half failures. What is the probability that this
              actually happens?
        \item What is the probability that $X$ is at least two standard
              deviations away from the mean?
        \item What is the probability that $X$ is at least three standard
              deviations away from the mean?
        \item How do these exact probabilities compare with the bounds given by
              Chebyshev's Theorem?
    \end{enumerate}
\end{exercise}

% ======================================================================
\begin{exercise}{8}
    Meteors arrive randomly and independently. The Eta Aquariids meteor shower
    has, at its peak and in idealised conditions, an average rate of $0.83$
    visible meteors per minute.
    \begin{enumerate}[label=(\alph*)]
        \item What is the distribution of the number of meteors seen in a given
              minute around the peak?
        \item What is the probability that no meteors are seen in a given minute?
        \item How long should an observer be prepared to wait to have
              probability $0.95$ of seeing a meteor?
    \end{enumerate}
\end{exercise}

\begin{exercise}{9}
    Each ball bowled to a Test cricket batter dismisses them with probability
    $1.25\%$, independently.
    \begin{enumerate}[label=(\alph*)]
        \item What is the distribution of the number of balls the batter faces
              until they are dismissed?
        \item What is the probability of a ``golden duck'' (dismissed by the
              first ball)?
        \item What is the expected length of their innings?
        \item Suppose the batter has faced 200 balls without being dismissed.
              How many further balls would you expect them to face?
    \end{enumerate}
\end{exercise}

% ======================================================================
\begin{exercise}{10}
    For each random variable, state which distribution is the best model:
    uniform, binomial, Poisson, geometric. If none fits, discuss why and suggest
    an appropriate distribution.
    \begin{enumerate}[label=(\alph*)]
        \item the face value of the top card in a well-shuffled standard deck
              (Ace, Jack, Queen, King $=1,11,12,13$);
        \item the number of days in a particular week on which your regular
              morning train is late (i.i.d.\ days; late = at least a minute after
              schedule);
        \item the number of cosmic rays that hit your laptop's CPU during your
              next class;
        \item the number of working days you wait, from tomorrow, until your
              train arrives on time (same assumptions as (b));
        \item the number of cereal packets you buy until you get one of your
              three favourite coupons (original coupon collector's scenario);
        \item the number of cowrie shells landing aperture-up when seven are
              thrown in Pachisi (identical, independent shells);
        \item a random digit from the first 1{,}002 digits after the decimal
              point of $3/7$;
        \item a random digit from the first 1{,}000 digits after the decimal
              point of $5/11$;
        \item a random digit from the first $10^{100}$ digits of $\pi$.
    \end{enumerate}
\end{exercise}

\begin{exercise}{11}
    (Same assumptions as the Birthday Paradox, Exercise 9.7.) You record
    birthdays of members of a customer loyalty scheme. How many people do you
    expect to have to enrol until their birthdays cover every day of the year?
\end{exercise}
